% 333_Kfrak_Dualitaet_En_ch.tex
% Doc. 333, FFGFT corpus, August 2026 — chapter file
% K_frak, Unrolling and the T~·m Duality
% Why the fractal correction is not a property of time unrolling

% ============================================================
\section*{Preliminary note}

In the FFGFT corpus, ``unrolling'' has hitherto always meant
\textbf{type-I decompactification}: the measurement act selects the
mass circle as time, $S^1_m\to\mathbb{R}_t$, discarding the winding
number — a lossy step (Doc.~285, Doc.~248, Doc.~270).
Doc.~332 introduces alongside this a different operation: the
\textbf{canonical function pullback}
$\Phi=p^*:L^2(S^1_m)\to\mathrm{Per}_{R_m}(\mathbb{R}_t)$, which is
lossless and exact — a pure change of coordinates with no loss of
information.

This document places both operations side by side, clarifies the role
of $K_{\text{frak}}$ in each case, and shows that the recursion
follows necessarily from the $\tilde{T}\cdot m$ duality.

% ============================================================
\section{Two Meanings of ``Unrolling'' in the Corpus}

\subsection{Type I: the measurement act (lossy)}

The meaning used so far in the corpus (Doc.~285 §\,``How FFGFT
unrolls''): the embedded observer selects the mass circle $S^1_m$ as
the time axis. The map
\[
  S^1_m \;\longrightarrow\; \mathbb{R}_t, \qquad
  \tilde T \;\mapsto\; t,
\]
is \textbf{lossy}: the winding number is discarded, the discrete
spectrum of the compact circle does not appear in $\mathbb{R}_t$.
This is not a mathematical operation of the theory but the bridge to
observation — the torus remains the fundamental structure,
$\mathbb{R}_t$ is the image (Doc.~285). At the same time the recursion
$\xi^n$ unrolls into a continuous dilation flow $t\mapsto\xi\,t$ with
log-period $|\ln\xi|=\ln 7500\approx 8.92$ (Doc.~285 §\,``The
recursion unrolls along with it'').

\subsection{Type III: the coordinate change (lossless)}

The meaning introduced in Doc.~332: the canonical function pullback
\[
  \Phi = p^* \;:\; L^2(S^1_m) \;\longrightarrow\;
  \mathrm{Per}_{R_m}(\mathbb{R}_t), \qquad
  (\Phi f)(t) = f(t \bmod R_m),
\]
is an \textbf{exact isometry} — no loss of information, no winding
number discarded, Fourier basis and Parseval identity preserved. It is
a change of representation within isomorphic Hilbert spaces (R96,
Doc.~332). The lossy direction lies at $\mathbb{R}_t\to S^1_m$, not
at the pullback itself.

\begin{center}
\begin{tabular}{lp{4.5cm}p{4.5cm}}
\toprule
 & \textbf{Type I (measurement act)} & \textbf{Type III (pullback)} \\
\midrule
Operation & $S^1_m\to\mathbb{R}_t$ & $L^2(S^1_m)\cong\mathrm{Per}_{R_m}(\mathbb{R}_t)$ \\
Loss & winding number discarded & none \\
Character & physical act & mathematical change of representation \\
Winding number & not preserved & preserved \\
Information loss & yes & no \\
Corpus references & Doc.~285, 248, 270 & Doc.~332 (R96) \\
\bottomrule
\end{tabular}
\end{center}

% ============================================================
\section{The Role of $K_{\text{frak}}$ in Both Cases}

\subsection{At the type-III pullback: invisible}

A Laplace operator on the compact $T^4$ is local — it cannot carry a
cumulative quantity (Doc.~314). $K_{\text{frak}}$ therefore does not
enter through the operator but through the \textbf{bridge constants}
$E_0$, $v$ at the scale conversion (Doc.~314, R72).

In the rolled-up state on $S^1_m$ the period $R_m=\hbar/(mc^2)$ is
exact in natural units — no $K_{\text{frak}}$. Only when $m$ is
substituted from the mass formula does $R_m$ carry the correction:
\[
  m_i = r_i\cdot\xi_0^{p_i}\cdot v\cdot K_{\text{frak}}
  \quad\Rightarrow\quad
  R_m^{\text{corr.}} = \frac{\hbar}{m_i^{\text{bare}}\cdot K_{\text{frak}}\cdot c^2}
  = \frac{R_m^{\text{bare}}}{K_{\text{frak}}}.
\]
$K_{\text{frak}}$ sits in the \emph{choice of $R_m$}, not in the
pullback mechanism. The pullback itself is exact in both cases (bare
or corrected).

\subsection{At the type-I measurement act: visible as a scale correction}

At the lossy unrolling, $K_{\text{frak}}$ appears as a correction
factor of the bridge constants: the ratio of the two $E_0$ values
(bare and $\alpha$-anchored) is $1/\sqrt{K_{\text{frak}}}$ (R72,
Doc.~314). The correction belongs to the \emph{scale structure} —
the transition between the compact and the observable representation
— not to the time geometry itself.

\begin{center}
\begin{tabular}{lp{4.5cm}p{4.5cm}}
\toprule
 & \textbf{Type I} & \textbf{Type III} \\
\midrule
$K_{\text{frak}}$ & visible in bridge constants
  ($E_0$ ratio, mass formula) & invisible at the isometry step;
  sits in $R_m$ via $m$ \\
Where exactly & at the SI transition & in the choice of period \\
\bottomrule
\end{tabular}
\end{center}

% ============================================================
\section{Mass and Time Unroll Together}

In the FFGFT compactification, $\tilde{T}$ and $m$ are a
\textbf{single geometric dimension}:
\[
  \tilde{T}\cdot m = 1 \qquad \text{(fundamental relation, Doc.~285)}.
\]
They are not two independent quantities. At the type-III pullback,
the pullback therefore unrolls $m$ and $t$ \emph{together}: the entire
mode $f\in L^2(S^1_m)$ is unrolled and the ratio $\tilde{T}/m$
remains constant. The physics does not change, only the coordinate
representation.

Unrolling $t$ alone — dropping $R_m$ — would violate
$\tilde{T}\cdot m=1$ and is structurally inadmissible.
$K_{\text{frak}}$ is \emph{not} the price of unrolling, but the price
of the \textbf{artificial separation} of $m$ and $t$ at the transition
to absolute SI units.

% ============================================================
\section{The Recursion Is Necessarily Bound to the Duality}

\subsection{Stable particle: Case B}

For a stable, isolated particle $\xi$ is frozen (Case~B, Doc.~315,
Doc.~313): the winding number is rational $74/75$, closure occurs after
exactly 75 revolutions. $K_{\text{frak}}=74/75$ is the closed
bookkeeping of this rolled-up geometry — exact in fractional
arithmetic, because $\xi$ is set rationally (Doc.~315).

\subsection{Dynamics: Case C and the recursion chain}

As soon as dynamics enter — interaction, RG flow, mass change — a
deeper consequence of $\tilde{T}\cdot m=1$ comes into play: if $m$
changes, $\tilde{T}=1/m$ changes necessarily with it. The period
$R_m=\hbar/(mc^2)$ changes, the time scale of the physics changes,
and this feeds back onto the effective mass. The recursion
(Doc.~295, Doc.~146) reads:
\[
  \xi_{n+1} = \xi_n\,(1-100\,\xi_n).
\]
The closure deficit per revolution no longer remains constant but runs
cumulatively. The rational 75-closure becomes a logarithmic spiral
(Case~C, Doc.~295):
\[
  \rho = 75\,e^{\beta/2\pi}
  \qquad\text{(equiangular spiral, self-similarity ratio }e\text{)}.
\]

In Case~C, $K_{\text{frak}}$ is no longer a constant but an
effectively scale-dependent quantity. The multiplicative form
$(1-\xi)^{100}$ is the natural bookkeeping of this running case; its
distance from $74/75$ — the binomial term $4950\,\xi^2$ (Doc.~315) —
is precisely the difference between frozen and running recursion.

\subsection{Why Case A is excluded}

Case~A ($K_{\text{frak}}=1$, closure after one revolution) would mean:
$\tilde{T}\cdot m=1$ generates no recursion, because $m$ never changes
— a static, interaction-free universe with $D_f=3$ exactly. This
contradicts the observed particle structure and the basic framework.

\subsection{Summary of the three cases}

\begin{center}
\begin{tabular}{p{1.2cm}p{3.5cm}p{3.5cm}p{3.5cm}}
\toprule
\textbf{Case} & \textbf{$\xi$} & \textbf{$K_{\text{frak}}$} &
\textbf{Binding to duality} \\
\midrule
A & constant, $d=0$ & $=1$ & no feedback, excluded \\
B & frozen & $=74/75$ exact & stable mode, duality at rest \\
C & running (recursion) & scale-dependent & duality generates running
  feedback $m\leftrightarrow\tilde{T}$ \\
\bottomrule
\end{tabular}
\end{center}

$K_{\text{frak}}=74/75$ is the \textbf{projection of Case~C onto
Case~B}: exact in the rolled-up, static description; first
approximation in the dynamic one.

% ============================================================
\section{Mode Specificity, Non-Closure and Superposition}

\subsection{Each mode carries its own period}

The period $R_m = \hbar/(mc^2)$ is mode-specific: every particle $i$
with torus mode $\mathcal{T}_i=(n_i,l_i,j_i)$ carries its own period.
Unrolling (Type~III or Type~I) is therefore always particle-specific
— it unrolls one particular mode, not a mixture. Simultaneously
unrolling different periods $R_{m_e}$ and $R_{m_\mu}$ would be
structurally inadmissible, because $\tilde{T}\cdot m=1$ would be
violated for both at once.

\subsection{The non-closure condition (Doc.~193)}

The same mode specificity underlies the non-closure condition
(Doc.~193): the prefactors $r_i$ are not numbers in a field but
labels of discrete geometric spaces — fully determined by
$\mathcal{T}_i$. \textbf{Multiplicative} combination of different
modes,
\[
  (r_i\cdot r_j,\; p_i+p_j) \;\notin\; \mathcal{P}
  \qquad\text{for }i\neq j,
\]
yields no FFGFT particle — the product lies outside the set of
admissible pairs \textbf{[B]}. $K_{\text{frak}}$ sits in the
mode-specific choice of $R_m$; a cross-mode $K_{\text{frak}}$ would
be the same inadmissible cross-product.

\subsection{Why superposition is nevertheless permitted — and where the difference lies}

In quantum mechanics, superpositions of different states are not
only permitted but fundamental. An electron in the sp$^3$ hybrid
orbital of a carbon atom is a linear combination of s and p states;
an excited atom is in a superposition of ground state and excited
state. The operation is \textbf{additive}:
\[
  \ket{\psi} = c_1\ket{\psi_1} + c_2\ket{\psi_2},
  \quad |c_1|^2 + |c_2|^2 = 1.
\]
This produces no new $(r,p)$-pair, no new mass, no new torus mode —
and therefore does not contradict the non-closure condition.

\textbf{The crucial difference between QM and FFGFT:}
In standard QM, time $t$ is not a quantum-mechanical operator and
not part of the Hilbert space — it is simply \emph{absent} as a
physical quantity of the system. The Schrödinger equation treats
$t$ as a classical external parameter that evolves the state from
outside, without itself belonging to the system. A superposition
$c_1\ket{\psi_1}+c_2\ket{\psi_2}$ therefore carries no
mode-bound time structure — $t$ does not belong to the
superposition but stands outside it.

In FFGFT, time is geometrically bound to mass: $\tilde{T}_i=1/m_i$.
Two modes with different masses $m_1\neq m_2$ carry different time
scales $\tilde{T}_1\neq\tilde{T}_2$. A superposition of such modes
would not merely mix amplitudes but overlay incompatible time scales
— structurally more problematic than in QM, where this conflict does
not arise.

In the FFGFT corpus, superposition is therefore described as
\textbf{intermediate torsion within a mode} (Doc.~174): a metastable
intermediate state between stable torsion extremes of the
\emph{same} mode, not a cross-mode superposition. Decoherence drives
this intermediate state back into one of the stable configurations.

\begin{center}
\begin{tabular}{lp{4.2cm}p{4.2cm}}
\toprule
 & \textbf{Standard QM} & \textbf{FFGFT} \\
\midrule
Time & absent: no operator, no Hilbert space element & $\tilde{T}=1/m$, mode-bound, geometric \\
Superposition & additive, $t$ stands outside & intermediate torsion
  within one mode \\
Cross-mode & algebraically permitted & incompatible $\tilde{T}$,
  structurally problematic \\
Non-closure & not relevant & $r_i\cdot r_j\notin\mathcal{P}$,
  multiplicatively forbidden \textbf{[B]} \\
\bottomrule
\end{tabular}
\end{center}

\subsection{The consistent boundary}

The linearity of quantum mechanics — the superposition principle —
and the non-closure condition protect the same boundary from two
sides: atoms and bound states generate no new fundamental particles,
no new torus modes, no new $(r,p)$-pair. The connection between
FFGFT and established atomic physics runs through the derived
quantities $\alpha$ and $m_e$ (Doc.~011, 044, 006) — quantum
mechanics takes these as input and derives binding energies and
spectra from them, without touching the mode structure.

\begin{quote}
\emph{``Unrolling'' in the FFGFT corpus has hitherto always meant the
lossy measurement act (Type~I, Doc.~285): winding number discarded,
$S^1_m\to\mathbb{R}_t$. The function pullback from Doc.~332 is a
different, lossless operation (Type~III): an exact change of
representation that unrolls $m$ and $t$ together and at which
$K_{\text{frak}}$ plays no role. $K_{\text{frak}}$ sits in the choice
of the period $R_m$ via the mass formula and becomes visible only at
the transition to absolute SI values. The recursion
$\xi_{n+1}=\xi_n(1-100\,\xi_n)$ is necessarily bound to
$\tilde{T}\cdot m=1$: every mass change entails a time change, which
feeds back recursively onto the mass. Case~B (stable particle,
$K_{\text{frak}}=74/75$) is the approximation for frozen duality;
Case~C (running recursion, log-spiral) is the complete dynamic
description.''}
\end{quote}

\medskip\noindent
\textbf{References:} Doc.~133 ($K_{\text{frak}}$ derivation),
Doc.~146 (recursion), Doc.~248 (observer),
Doc.~270 (type-I theorem), Doc.~285 (decompactification),
Doc.~295 (log-spiral), Doc.~313 (three cases),
Doc.~314 (local operator, bridge constants),
Doc.~315 (form question $K_{\text{frak}}$),
Doc.~332 (pullback isometry, R96).
